\documentclass[12pt]{article}
\usepackage{amsmath, amssymb, amsthm}
\usepackage{mathtools}
\usepackage{enumerate}
\usepackage{graphicx}
\usepackage{subcaption} % only if you want them side-by-side
\usepackage{float}      % for [H] placement (optional)
\usepackage{geometry}
\geometry{margin=1in}

\begin{document}

\title{CS 395T Homework 4}
\author{Surya Sunkari (svs879)}
\date{April 4, 2026}

\maketitle

\section*{Problem 1}

\subsection*{(i)}

\begin{proof}
By the equivalent characterization of total variation (Remark 1, Lecture 14),
\[
D_{\mathrm{TV}}(\mathcal{T}\pi, \pi^\star)
= \sup_{A \subseteq \mathbb{R}^d} \left( \int_{A} (\mathcal{T}\pi)(\mathbf{y}) \,\mathrm{d}\mathbf{y} - \int_{A} \pi^\star(\mathbf{y}) \,\mathrm{d}\mathbf{y} \right).
\]
Fix a measurable set $A \subseteq \mathbb{R}^d$. Expanding $(\mathcal{T}\pi)(\mathbf{y})$ and using stationarity of $\pi^\star$:
\begin{align*}
\int_{A} (\mathcal{T}\pi)(\mathbf{y}) \,\mathrm{d}\mathbf{y} - \int_{A} \pi^\star(\mathbf{y}) \,\mathrm{d}\mathbf{y}
&= \int_{A} \int_{\mathbb{R}^d} \mathcal{T}_{\mathbf{x}}(\mathbf{y})\, \pi(\mathbf{x}) \,\mathrm{d}\mathbf{x} \,\mathrm{d}\mathbf{y} - \int_{A} \int_{\mathbb{R}^d} \mathcal{T}_{\mathbf{x}}(\mathbf{y})\, \pi^\star(\mathbf{x}) \,\mathrm{d}\mathbf{x} \,\mathrm{d}\mathbf{y} \\
&= \int_{\mathbb{R}^d} \left( \int_{A} \mathcal{T}_{\mathbf{x}}(\mathbf{y}) \,\mathrm{d}\mathbf{y} \right) \bigl(\pi(\mathbf{x}) - \pi^\star(\mathbf{x})\bigr) \,\mathrm{d}\mathbf{x},
\end{align*}
where the interchange of integrals is justified by Fubini's theorem (the integrands are non-negative densities, hence measurable, and the product measures are $\sigma$-finite).

Since $\mathcal{T}_{\mathbf{x}}$ is a probability distribution for each $\mathbf{x}$, we have $\int_{A} \mathcal{T}_{\mathbf{x}}(\mathbf{y}) \,\mathrm{d}\mathbf{y} \leq 1$. Define $g(\mathbf{x}) := \int_{A} \mathcal{T}_{\mathbf{x}}(\mathbf{y}) \,\mathrm{d}\mathbf{y} \in [0, 1]$. Then
\begin{align*}
\int_{\mathbb{R}^d} g(\mathbf{x}) \bigl(\pi(\mathbf{x}) - \pi^\star(\mathbf{x})\bigr) \,\mathrm{d}\mathbf{x}
&= \int_{\mathbb{R}^d} g(\mathbf{x}) \bigl(\pi(\mathbf{x}) - \pi^\star(\mathbf{x})\bigr)^+ \,\mathrm{d}\mathbf{x} - \int_{\mathbb{R}^d} g(\mathbf{x}) \bigl(\pi(\mathbf{x}) - \pi^\star(\mathbf{x})\bigr)^- \,\mathrm{d}\mathbf{x} \\
&\leq \int_{\mathbb{R}^d} \bigl(\pi(\mathbf{x}) - \pi^\star(\mathbf{x})\bigr)^+ \,\mathrm{d}\mathbf{x}, \quad \text{(since } g \leq 1 \text{ and } g \geq 0\text{)}
\end{align*}
where $(a)^+ = \max(a, 0)$ and $(a)^- = \max(-a, 0)$.

By the standard identity relating the positive part to TV distance,
\[
\int_{\mathbb{R}^d} \bigl(\pi(\mathbf{x}) - \pi^\star(\mathbf{x})\bigr)^+ \,\mathrm{d}\mathbf{x}
= \sup_{B \subseteq \mathbb{R}^d} \left( \int_{B} \pi(\mathbf{x}) \,\mathrm{d}\mathbf{x} - \int_{B} \pi^\star(\mathbf{x}) \,\mathrm{d}\mathbf{x} \right)
= D_{\mathrm{TV}}(\pi, \pi^\star),
\]
where the first equality holds by taking $B = \{\mathbf{x} : \pi(\mathbf{x}) > \pi^\star(\mathbf{x})\}$, and the second is the equivalent characterization (Remark 1, Lecture 14).

Since the bound $D_{\mathrm{TV}}(\pi, \pi^\star)$ holds for every measurable $A$, taking the supremum over $A$ gives
\[
D_{\mathrm{TV}}(\mathcal{T}\pi, \pi^\star) \leq D_{\mathrm{TV}}(\pi, \pi^\star). \qedhere
\]
\end{proof}

\subsection*{(ii)}

\begin{proof}
We show $(\mathcal{T}\pi)(\mathbf{y}) \leq \beta \, \pi^\star(\mathbf{y})$ for all $\mathbf{y} \in \mathbb{R}^d$. Since $\pi$ is $\beta$-warm with respect to $\pi^\star$ (Definition 2, Lecture 15), we have $\pi(\mathbf{x}) \leq \beta \, \pi^\star(\mathbf{x})$ for all $\mathbf{x} \in \mathbb{R}^d$. Then for any $\mathbf{y} \in \mathbb{R}^d$,
\begin{align*}
(\mathcal{T}\pi)(\mathbf{y})
&= \int_{\mathbb{R}^d} \mathcal{T}_{\mathbf{x}}(\mathbf{y})\, \pi(\mathbf{x}) \,\mathrm{d}\mathbf{x} \\
&\leq \int_{\mathbb{R}^d} \mathcal{T}_{\mathbf{x}}(\mathbf{y})\, \beta \, \pi^\star(\mathbf{x}) \,\mathrm{d}\mathbf{x} \quad \text{($\pi(\mathbf{x}) \leq \beta\, \pi^\star(\mathbf{x})$ and } \mathcal{T}_{\mathbf{x}}(\mathbf{y}) \geq 0\text{)} \\
&= \beta \int_{\mathbb{R}^d} \mathcal{T}_{\mathbf{x}}(\mathbf{y})\, \pi^\star(\mathbf{x}) \,\mathrm{d}\mathbf{x} \\
&= \beta \, (\mathcal{T}\pi^\star)(\mathbf{y}) \\
&= \beta \, \pi^\star(\mathbf{y}), \quad \text{(stationarity of $\pi^\star$)}
\end{align*}
so $\mathcal{T}\pi$ is $\beta$-warm with respect to $\pi^\star$.
\end{proof}

\section*{Problem 2}
\subsection*{(i)}

\textbf{Algorithm (Rejection Sampling):}
\begin{enumerate}
\item Sample $\omega \sim Q$ using one query to $\mathcal{O}$.
\item Query the value oracles for $f(\omega)$ and $g(\omega)$.
\item Accept $\omega$ with probability $\frac{f(\omega)}{C \cdot g(\omega)} \in [0,1]$.
\item If rejected, return to step 1. If accepted, output $\omega$.
\end{enumerate}

The acceptance probability is well-defined since $\frac{f(\omega)}{g(\omega)} \leq C$ implies $\frac{f(\omega)}{C \cdot g(\omega)} \leq 1$, and $f \geq 0$, $g > 0$ ensures $\frac{f(\omega)}{C \cdot g(\omega)} \geq 0$.

\textbf{Correctness:}
\begin{proof}
Let $Z_f = \int_\Omega f(\omega)\,\mathrm{d}\omega$ and $Z_g = \int_\Omega g(\omega)\,\mathrm{d}\omega$, so $P(\omega) = \frac{f(\omega)}{Z_f}$ and $Q(\omega) = \frac{g(\omega)}{Z_g}$. The probability that a sample $\omega \sim Q$ is accepted and falls in a measurable set $A \subseteq \Omega$ is:
\[
\int_A Q(\omega) \cdot \frac{f(\omega)}{C \cdot g(\omega)}\,\mathrm{d}\omega = \int_A \frac{g(\omega)}{Z_g} \cdot \frac{f(\omega)}{C \cdot g(\omega)}\,\mathrm{d}\omega = \frac{1}{C \cdot Z_g} \int_A f(\omega)\,\mathrm{d}\omega.
\]
Conditioning on acceptance, the density of the accepted sample over $A$ is:
\[
\frac{\frac{1}{C \cdot Z_g} \int_A f(\omega)\,\mathrm{d}\omega}{\frac{1}{C \cdot Z_g} \int_\Omega f(\omega)\,\mathrm{d}\omega} = \frac{\int_A f(\omega)\,\mathrm{d}\omega}{Z_f},
\]
which is precisely $\int_A P(\omega)\,\mathrm{d}\omega$. Hence the accepted sample is distributed according to $P$.
\end{proof}

\textbf{Query complexity:}
\begin{proof}
The acceptance probability in each iteration is:
\[
\Pr[\text{accept}] = \int_\Omega Q(\omega) \cdot \frac{f(\omega)}{C \cdot g(\omega)}\,\mathrm{d}\omega = \frac{1}{C \cdot Z_g} \int_\Omega f(\omega)\,\mathrm{d}\omega = \frac{Z_f}{C \cdot Z_g}.
\]
The number of iterations until acceptance is geometric with success probability $\frac{Z_f}{C \cdot Z_g}$, so the expected number of iterations is:
\[
\frac{C \cdot Z_g}{Z_f} = \frac{C \int_\Omega g(\omega)\,\mathrm{d}\omega}{\int_\Omega f(\omega)\,\mathrm{d}\omega} = N.
\]
Each iteration uses one query to $\mathcal{O}$ and one value oracle query to each of $f$ and $g$, giving at most $N$ queries of each type in expectation.
\end{proof}

\subsection*{(ii)}

We apply the rejection sampling algorithm from part (i) with the following choices:
\begin{itemize}
\item $\Omega = \mathbb{B}(\mathbf{0}_d, 1)$,
\item the unnormalized target density $\tilde{\pi}(\mathbf{x}) = \exp(-f(\mathbf{x}))$ (playing the role of $f$ in part (i)),
\item $g(\mathbf{x}) = 1$ for all $\mathbf{x} \in \Omega$, so $Q = \mathrm{Uniform}(\mathbb{B}(\mathbf{0}_d, 1))$.
\end{itemize}

\textbf{Computing $C$:} Since $f$ is $1$-Lipschitz and $\mathbb{B}(\mathbf{0}_d, 1)$ has diameter at most $2$, for all $\mathbf{x} \in \mathbb{B}(\mathbf{0}_d, 1)$:
\[
|f(\mathbf{x}) - f(\mathbf{0})| \leq \|\mathbf{x} - \mathbf{0}\|_2 \leq 1,
\]
so $f(\mathbf{x}) \geq f(\mathbf{0}) - 1$, which gives:
\[
\exp(-f(\mathbf{x})) \leq \exp(-f(\mathbf{0}) + 1) = e \cdot \exp(-f(\mathbf{0})).
\]
Since $g(\mathbf{x}) = 1$, we need $\frac{\exp(-f(\mathbf{x}))}{g(\mathbf{x})} \leq C$. Set
\[
C = e \cdot \exp(-f(\mathbf{0})),
\]
which requires one value oracle query to $f$ (at $\mathbf{0}$).

\textbf{Algorithm:}
\begin{enumerate}
\item Query $f(\mathbf{0})$ and set $C = e \cdot \exp(-f(\mathbf{0}))$.
\item Sample $\mathbf{x}$ uniformly from $\mathbb{B}(\mathbf{0}_d, 1)$.
\item Query $f(\mathbf{x})$ and accept with probability
\[
\frac{\exp(-f(\mathbf{x}))}{C \cdot g(\mathbf{x})} = \frac{\exp(-f(\mathbf{x}))}{e \cdot \exp(-f(\mathbf{0}))} = \exp(f(\mathbf{0}) - f(\mathbf{x}) - 1).
\]
This lies in $[0,1]$: the bound $\frac{\exp(-f(\mathbf{x}))}{g(\mathbf{x})} \leq C$ gives $\frac{\exp(-f(\mathbf{x}))}{C \cdot g(\mathbf{x})} \leq 1$, and $\exp(-f(\mathbf{x})) > 0$ gives $\geq 0$.

\item If rejected, return to step 2. If accepted, output $\mathbf{x}$.
\end{enumerate}

\textbf{Correctness} follows from part (i): the output is distributed according to $\pi \propto \exp(-f)$ over $\mathbb{B}(\mathbf{0}_d, 1)$.

\textbf{Query complexity:} By part (i), the expected number of iterations is:
\[
N = \frac{C \int_{\mathbb{B}(\mathbf{0}_d,1)} g(\mathbf{x})\,\mathrm{d}\mathbf{x}}{\int_{\mathbb{B}(\mathbf{0}_d,1)} \exp(-f(\mathbf{x}))\,\mathrm{d}\mathbf{x}} = \frac{e \cdot \exp(-f(\mathbf{0})) \cdot \mathrm{Vol}(\mathbb{B}(\mathbf{0}_d,1))}{\int_{\mathbb{B}(\mathbf{0}_d,1)} \exp(-f(\mathbf{x}))\,\mathrm{d}\mathbf{x}}.
\]
Since $f(\mathbf{x}) \leq f(\mathbf{0}) + 1$ for all $\mathbf{x} \in \mathbb{B}(\mathbf{0}_d,1)$, we have $\exp(-f(\mathbf{x})) \geq \exp(-f(\mathbf{0}) - 1)$, so:
\[
\int_{\mathbb{B}(\mathbf{0}_d,1)} \exp(-f(\mathbf{x}))\,\mathrm{d}\mathbf{x} \geq \exp(-f(\mathbf{0}) - 1) \cdot \mathrm{Vol}(\mathbb{B}(\mathbf{0}_d,1)).
\]
Therefore:
\[
N \leq \frac{e \cdot \exp(-f(\mathbf{0})) \cdot \mathrm{Vol}(\mathbb{B}(\mathbf{0}_d,1))}{\exp(-f(\mathbf{0}) - 1) \cdot \mathrm{Vol}(\mathbb{B}(\mathbf{0}_d,1))} = \frac{e \cdot \exp(-f(\mathbf{0}))}{\exp(-f(\mathbf{0}) - 1)} = e \cdot e = e^2.
\]
Including the initial query to compute $f(\mathbf{0})$, the total expected number of value oracle queries to $f$ is $1 + e^2 = O(1)$.

\section*{Problem 3}

\subsection*{(i)}
The kernel of $\mathbf{L}_G$ is the span of the indicator vectors of the connected components of $G$.

\begin{proof}
By Lemma 10 (Part II),
\[
\mathbf{x}^\top \mathbf{L}_G \mathbf{x} = \sum_{e = (u,v) \in E} w_e (x_u - x_v)^2.
\]
Since $w_e > 0$ for all $e \in E$, this equals zero if and only if $x_u = x_v$ for every edge $(u, v) \in E$. Equivalently, $\mathbf{x}$ is constant on each connected component of $G$. The set of all such vectors is precisely $\operatorname{span}\{\mathbf{1}_{C_1}, \ldots, \mathbf{1}_{C_k}\}$, where $C_1, \ldots, C_k$ are the connected components of $G$.
\end{proof}

\subsection*{(ii)}
For $\mathbf{x} \in \{-\tfrac{1}{2}, \tfrac{1}{2}\}^V$, the vector $\mathbf{x}$ defines a cut $(S, V \setminus S)$ where $S = \{v \in V : x_v = \tfrac{1}{2}\}$. For each edge $e = (u,v)$, we have $(x_u - x_v)^2 = 1$ if $e$ crosses the cut (i.e., exactly one of $u, v$ is in $S$) and $(x_u - x_v)^2 = 0$ otherwise. By Lemma 10 (Part II),
\[
\mathbf{x}^\top \mathbf{L}_G \mathbf{x} = \sum_{e = (u,v) \in E} w_e (x_u - x_v)^2 = \sum_{\substack{e = (u,v) \in E \\ |\{u,v\} \cap S| = 1}} w_e.
\]
This is the total weight of edges crossing the cut $(S, V \setminus S)$.

\subsection*{(iii)}
\begin{proof}
Write the Laplacian as a sum of rank-one terms:
\[
\mathbf{L}_G = \sum_{e = (u,v) \in E} w_e (\mathbf{e}_u - \mathbf{e}_v)(\mathbf{e}_u - \mathbf{e}_v)^\top.
\]
Define the matrix $\mathbf{A} \in \mathbb{R}^{|E| \times |V|}$ whose row indexed by edge $e = (u,v)$ is $\mathbf{a}_e := \sqrt{w_e}(\mathbf{e}_u - \mathbf{e}_v)^\top$. Then
\[
\mathbf{A}^\top \mathbf{A} = \sum_{e \in E} \mathbf{a}_e^\top \mathbf{a}_e = \sum_{e = (u,v) \in E} w_e (\mathbf{e}_u - \mathbf{e}_v)(\mathbf{e}_u - \mathbf{e}_v)^\top = \mathbf{L}_G.
\]
We apply Proposition 3 (Part IX) to $\mathbf{A}$. Set $\varepsilon = 0.1$, $\delta = \tfrac{1}{2}$, $n = |E|$, and $d = |V|$.

\textbf{Rank bound:} $\operatorname{rank}(\mathbf{A}) = \operatorname{rank}(\mathbf{L}_G) \leq |V|$.

\textbf{Sample complexity:} By Proposition 3 (Part IX), taking
\[
K \geq \frac{3 \operatorname{rank}(\mathbf{A})}{\varepsilon^2} \log\!\left(\frac{2d}{\delta}\right) \leq \frac{3|V|}{0.01} \log(4|V|) = O(|V| \log |V|)
\]
i.i.d.\ samples $e_1, \ldots, e_K$ from the leverage score distribution $p_e = \tau_e(\mathbf{A}) / \|\boldsymbol{\tau}\|_1$ guarantees, with probability at least $\tfrac{1}{2}$, that
\[
(1 - \varepsilon)\,\mathbf{A}^\top \mathbf{A} \preceq \mathbf{P} \preceq (1 + \varepsilon)\,\mathbf{A}^\top \mathbf{A},
\]
where
\[
\mathbf{P} = \sum_{k=1}^{K} \frac{1}{K\, p_{e_k}} \mathbf{a}_{e_k}^\top \mathbf{a}_{e_k} = \sum_{k=1}^{K} \frac{w_{e_k}}{K\, p_{e_k}} (\mathbf{e}_{u_k} - \mathbf{e}_{v_k})(\mathbf{e}_{u_k} - \mathbf{e}_{v_k})^\top.
\]

\textbf{Identifying the sparsifier:} Each term in $\mathbf{P}$ has the form $w'(\mathbf{e}_u - \mathbf{e}_v)(\mathbf{e}_u - \mathbf{e}_v)^\top$ for some edge $(u,v) \in E$ with positive weight $w' = w_{e_k}/(K\, p_{e_k}) > 0$. If the same edge is sampled multiple times, we sum its weights. Hence $\mathbf{P} = \mathbf{L}_H$ for a reweighted subgraph $H = (V, E', \mathbf{w}')$ with $E' \subseteq E$ and $\mathbf{w}' \in \mathbb{R}_{>0}^{E'}$. The number of distinct edges satisfies $|E'| \leq K = O(|V| \log |V|)$.

Since $\mathbf{A}^\top \mathbf{A} = \mathbf{L}_G$ and $\varepsilon = 0.1$, the guarantee becomes
\[
0.9\, \mathbf{L}_G \preceq \mathbf{L}_H \preceq 1.1\, \mathbf{L}_G.
\]
The sampling succeeds with probability at least $\tfrac{1}{2} > 0$, so such a graph $H$ exists.
\end{proof}

\section*{Problem 4}

Set $\pi_0 = \mathcal{N}\!\left(\mathbf{x}^\star,\, \frac{1}{L}\mathbf{I}_d\right)$, the Gaussian centered at the minimizer with covariance $\frac{1}{L}\mathbf{I}_d$. We show $\pi_0$ is $\beta$-warm with respect to $\pi^\star$ for $\beta = \kappa^{d/2} = \exp\!\bigl(\tfrac{d}{2}\log\kappa\bigr)$.

\begin{proof}
Write the normalizing constant $Z := \int_{\mathbb{R}^d} \exp(-f(\mathbf{x}))\,d\mathbf{x}$, so that
\[
\pi^\star(\mathbf{x}) = \frac{\exp(-f(\mathbf{x}))}{Z}.
\]
The Gaussian density is
\[
\pi_0(\mathbf{x}) = \left(\frac{L}{2\pi}\right)^{d/2} \exp\!\left(-\frac{L}{2}\|\mathbf{x} - \mathbf{x}^\star\|_2^2\right).
\]
The warmness ratio at any $\mathbf{x} \in \mathbb{R}^d$ is
\[
\frac{\pi_0(\mathbf{x})}{\pi^\star(\mathbf{x})} = \left(\frac{L}{2\pi}\right)^{d/2} \exp\!\left(-\frac{L}{2}\|\mathbf{x} - \mathbf{x}^\star\|_2^2 + f(\mathbf{x})\right) \cdot Z.
\]

\textbf{Bounding the exponent.} Since $f$ is $L$-smooth and $\mathbf{x}^\star$ is the minimizer (so $\nabla f(\mathbf{x}^\star) = \mathbf{0}$), the quadratic upper bound gives
\[
f(\mathbf{x}) \leq f(\mathbf{x}^\star) + \frac{L}{2}\|\mathbf{x} - \mathbf{x}^\star\|_2^2.
\]
Therefore, for every $\mathbf{x} \in \mathbb{R}^d$,
\[
-\frac{L}{2}\|\mathbf{x} - \mathbf{x}^\star\|_2^2 + f(\mathbf{x}) \leq f(\mathbf{x}^\star).
\]

\textbf{Bounding $Z$.} Since $f$ is $\mu$-strongly convex with minimizer $\mathbf{x}^\star$ (so $\nabla f(\mathbf{x}^\star) = \mathbf{0}$),
\[
f(\mathbf{x}) \geq f(\mathbf{x}^\star) + \frac{\mu}{2}\|\mathbf{x} - \mathbf{x}^\star\|_2^2 \quad \text{for all } \mathbf{x} \in \mathbb{R}^d.
\]
Exponentiating and integrating:
\begin{align*}
Z &= \int_{\mathbb{R}^d} \exp(-f(\mathbf{x}))\,d\mathbf{x} \\
  &\leq \int_{\mathbb{R}^d} \exp\!\left(-f(\mathbf{x}^\star) - \frac{\mu}{2}\|\mathbf{x} - \mathbf{x}^\star\|_2^2\right)d\mathbf{x} \\
  &= \exp(-f(\mathbf{x}^\star)) \int_{\mathbb{R}^d} \exp\!\left(-\frac{\mu}{2}\|\mathbf{x} - \mathbf{x}^\star\|_2^2\right)d\mathbf{x} \\
  &= \exp(-f(\mathbf{x}^\star)) \left(\frac{2\pi}{\mu}\right)^{d/2}.
\end{align*}

\textbf{Combining.} Substituting both bounds into the ratio:
\begin{align*}
\frac{\pi_0(\mathbf{x})}{\pi^\star(\mathbf{x})}
&= \left(\frac{L}{2\pi}\right)^{d/2} \exp\!\left(-\frac{L}{2}\|\mathbf{x} - \mathbf{x}^\star\|_2^2 + f(\mathbf{x})\right) \cdot Z \\
&\leq \left(\frac{L}{2\pi}\right)^{d/2} \exp(f(\mathbf{x}^\star)) \cdot \exp(-f(\mathbf{x}^\star))\left(\frac{2\pi}{\mu}\right)^{d/2} \\
&= \left(\frac{L}{2\pi}\right)^{d/2} \cdot \left(\frac{2\pi}{\mu}\right)^{d/2} \\
&= \left(\frac{L}{\mu}\right)^{d/2} = \kappa^{d/2}.
\end{align*}
This holds for every $\mathbf{x} \in \mathbb{R}^d$, so by Definition 2 (Part XV), $\pi_0$ is $\beta$-warm with respect to $\pi^\star$ with
\[
\beta = \kappa^{d/2} = \exp\!\left(\frac{d}{2}\log\kappa\right) = \exp(O(d\log\kappa)).
\]
\end{proof}

\section*{Problem 5}
\begin{proof}
Set $\pi_k := T^k \pi_0$ and $\delta_k := D_{\mathrm{TV}}(\pi_k, \pi^\star)$.

By Lemma 2 (Part XV), since $\pi_0$ is $\beta$-warm with respect to $\pi^\star$, $\pi_k$ is $\beta$-warm with respect to $\pi^\star$ for all $k \ge 0$ (apply the lemma inductively: $T\pi_k$ is $\beta$-warm whenever $\pi_k$ is).

\textbf{Claim:} If $\delta_k \ge \epsilon$, then $\delta_{k+1} \le \delta_k / 2$.

Since $D_{\mathrm{TV}}(\pi_k, \pi^\star) = \delta_k$, decompose via the TV coupling:
\[
\pi_k = (1 - \delta_k)\sigma_k + \delta_k \rho_k, \qquad \pi^\star = (1 - \delta_k)\sigma_k + \delta_k \rho_k',
\]
where
\[
\sigma_k(x) = \frac{\min(\pi_k(x), \pi^\star(x))}{1 - \delta_k}, \quad \rho_k(x) = \frac{(\pi_k(x) - \pi^\star(x))^+}{\delta_k}, \quad \rho_k'(x) = \frac{(\pi^\star(x) - \pi_k(x))^+}{\delta_k}.
\]
These are valid densities: $\sigma_k, \rho_k, \rho_k' \ge 0$ by construction, and integrating the decomposition of $\pi_k$ gives
\[
1 = (1 - \delta_k)\int \sigma_k + \delta_k \int \rho_k,
\]
while $\int \sigma_k = \frac{\int \min(\pi_k, \pi^\star)}{1 - \delta_k} = \frac{1 - \delta_k}{1 - \delta_k} = 1$ (since $\int \min(\pi_k, \pi^\star) = 1 - D_{\mathrm{TV}}(\pi_k, \pi^\star) = 1 - \delta_k$), so $\int \rho_k = 1$. The same argument gives $\int \rho_k' = 1$.

By linearity of the transition operator $T$ and stationarity $T\pi^\star = \pi^\star$:
\begin{align*}
T\pi_k &= (1 - \delta_k) T\sigma_k + \delta_k T\rho_k, \\
\pi^\star = T\pi^\star &= (1 - \delta_k) T\sigma_k + \delta_k T\rho_k'.
\end{align*}
Subtracting, for any measurable $S$:
\[
T\pi_k(S) - \pi^\star(S) = \delta_k \bigl(T\rho_k(S) - T\rho_k'(S)\bigr).
\]
Taking $\sup_S$ on both sides:
\[
\delta_{k+1} = D_{\mathrm{TV}}(T\pi_k, \pi^\star) = \delta_k \cdot D_{\mathrm{TV}}(T\rho_k, T\rho_k').
\]

For $\rho_k$:
\[
\rho_k(x) = \frac{(\pi_k(x) - \pi^\star(x))^+}{\delta_k} \le \frac{\pi_k(x)}{\delta_k} \le \frac{\beta \pi^\star(x)}{\delta_k},
\]
so $\rho_k$ is $\frac{\beta}{\delta_k}$-warm with respect to $\pi^\star$. Since $\delta_k \ge \epsilon$, $\rho_k$ is $\frac{\beta}{\epsilon}$-warm. For $\rho_k'$:
\[
\rho_k'(x) = \frac{(\pi^\star(x) - \pi_k(x))^+}{\delta_k} \le \frac{\pi^\star(x)}{\delta_k},
\]
so $\rho_k'$ is $\frac{1}{\delta_k}$-warm, hence $\frac{\beta}{\epsilon}$-warm since $\beta \ge 1$ and $\delta_k \ge \epsilon$.

By the hypothesis on the random walk, any $\frac{\beta}{\epsilon}$-warm density $\pi$ satisfies $D_{\mathrm{TV}}(T\pi, \pi^\star) \le \frac{1}{4}$. Applying this to $\rho_k$ and $\rho_k'$:
\[
D_{\mathrm{TV}}(T\rho_k, \pi^\star) \le \frac{1}{4}, \qquad D_{\mathrm{TV}}(T\rho_k', \pi^\star) \le \frac{1}{4}.
\]
By the triangle inequality:
\[
D_{\mathrm{TV}}(T\rho_k, T\rho_k') \le D_{\mathrm{TV}}(T\rho_k, \pi^\star) + D_{\mathrm{TV}}(\pi^\star, T\rho_k') \le \frac{1}{4} + \frac{1}{4} = \frac{1}{2}.
\]
Therefore $\delta_{k+1} = \delta_k \cdot D_{\mathrm{TV}}(T\rho_k, T\rho_k') \le \delta_k / 2$, proving the claim.

We show by induction that $\delta_k \le 2^{-(k+1)}$ for all $k \ge 1$. The base case $k = 1$: since $\pi_0$ is $\beta$-warm and $\epsilon < 1$ implies $\frac{\beta}{\epsilon} \ge \beta \ge 1$, $\pi_0$ is $\frac{\beta}{\epsilon}$-warm, so the hypothesis gives $\delta_1 \le \frac{1}{4} = 2^{-2}$. For the inductive step, suppose $\delta_k \le 2^{-(k+1)}$. If $\delta_k < \epsilon$, then $\delta_{k+1} \le \delta_k < \epsilon$ by Lemma 2 (Part XV) (TV distance does not increase under $T$), so the conclusion holds. If $\delta_k \ge \epsilon$, the halving gives $\delta_{k+1} \le \delta_k / 2 \le 2^{-(k+2)}$.

For $k \ge \log_2(1/\epsilon)$:
\[
\delta_k \le 2^{-(k+1)} \le 2^{-(\log_2(1/\epsilon) + 1)} = \frac{\epsilon}{2} \le \epsilon.
\]
Hence $D_{\mathrm{TV}}(T^k \pi_0, \pi^\star) \le \epsilon$ whenever $k \ge \log_2(1/\epsilon)$.
\end{proof}

\end{document}
